%% paper81_hadwiger_conjecture_en.tex
%% Hadwiger's Conjecture for n >= 7: Graph Coloring and Minor Theory
%% Author: Michael Fuhrmann
%% Specialist Mathematics Project, Build 166
%% Date: 2026-03-12

\documentclass[12pt,a4paper]{amsart}
\usepackage{amsmath,amssymb,amsthm,mathtools,enumitem,booktabs,array,hyperref}
\usepackage[utf8]{inputenc}
\usepackage[T1]{fontenc}
\usepackage{geometry}
\geometry{margin=2.5cm}

%% Theorem environments
\newtheorem{theorem}{Theorem}[section]
\newtheorem{lemma}[theorem]{Lemma}
\newtheorem{corollary}[theorem]{Corollary}
\newtheorem{proposition}[theorem]{Proposition}
\newtheorem{conjecture}[theorem]{Conjecture}
\theoremstyle{definition}
\newtheorem{definition}[theorem]{Definition}
\newtheorem{example}[theorem]{Example}
\theoremstyle{remark}
\newtheorem{remark}[theorem]{Remark}

\hypersetup{
  colorlinks=true,
  linkcolor=blue,
  citecolor=blue,
  urlcolor=blue
}

\begin{document}

\title{Hadwiger's Conjecture for $k \geq 7$:\\
Graph Coloring, Complete Minors,\\
and the Four-Color Theorem Connection}

\author{Michael Fuhrmann}
\address{Specialist Mathematics Project, Build 166}
\email{specialist-maths@project.local}
\date{2026-03-12}

\begin{abstract}
Hadwiger's conjecture (1943) asserts that every graph $G$ with chromatic number
$\chi(G) \geq k$ contains the complete graph $K_k$ as a minor. This is one of
the deepest open problems in graph theory, combining the theory of graph colorings
with the structure theory of graph minors developed by Robertson and Seymour.

The conjecture has been verified for $k \leq 6$: the cases $k \leq 4$ follow
from classical results, $k = 5$ is equivalent to the Four-Color Theorem (Wagner 1937),
and $k = 6$ was proved by Robertson, Seymour, and Thomas (1993) also using the
Four-Color Theorem. For $k \geq 7$, the conjecture remains entirely open.

We present complete proofs for $k \leq 4$, a detailed analysis of the $k = 5$
equivalence with the Four-Color Theorem, the Robertson--Seymour--Thomas proof
for $k = 6$, and a survey of partial results for large $k$ including
Reed--Seymour average-degree bounds and Kühn--Osthus results on graphs with
large chromatic number.
\end{abstract}

\maketitle

\tableofcontents

%% =============================================================================
\section{Introduction}
%% =============================================================================

In 1943, Hugo Hadwiger~\cite{Hadwiger1943} proposed the following unifying
conjecture connecting the chromatic number of a graph with its minor structure:

\begin{conjecture}[Hadwiger, 1943]\label{conj:hadwiger}
For every integer $k \geq 1$, if $\chi(G) \geq k$ then $G$ contains $K_k$ as
a minor (written $G \succeq K_k$).
\end{conjecture}

Equivalently: every graph without a $K_k$ minor is $(k-1)$-colorable.

\begin{definition}
Let $G = (V,E)$ be a graph.
\begin{enumerate}[label=(\alph*)]
  \item The \emph{chromatic number} $\chi(G)$ is the smallest number of colors
        needed to color the vertices of $G$ so that adjacent vertices receive
        different colors.
  \item A graph $H$ is a \emph{minor} of $G$ (written $H \preceq G$) if $H$ can
        be obtained from $G$ by a sequence of edge contractions, edge deletions,
        and vertex deletions.
  \item The \emph{Hadwiger number} $\eta(G)$ (also written $h(G)$) is the largest
        $k$ such that $K_k \preceq G$.
\end{enumerate}
\end{definition}

Hadwiger's conjecture can be restated as: $\chi(G) \leq \eta(G)$ for all graphs $G$.

Note that $\chi(G) \geq \omega(G)$ (clique number) and $\eta(G) \geq \omega(G)$
trivially, so the conjunction is not tautological.

%% =============================================================================
\section{Graph Minors: Basic Theory}
%% =============================================================================

\begin{definition}
An \emph{edge contraction} contracts an edge $uv$ by merging $u$ and $v$ into
a single vertex $w$ whose neighbors are $N(u) \cup N(v) \setminus \{u,v\}$.
\end{definition}

\begin{proposition}[Properties of minors]\label{prop:minorprop}
\begin{enumerate}[label=(\roman*)]
  \item The minor relation $\preceq$ is a partial order on isomorphism classes
        of graphs.
  \item If $H \preceq G$ and $\chi(H) \geq k$, then $\chi(G) \geq k$.
  \item Equivalently, $K_k \preceq G$ implies $\chi(G) \geq k$ is false in general;
        Hadwiger's conjecture is the reverse direction.
\end{enumerate}
\end{proposition}

\begin{theorem}[Robertson--Seymour Graph Minor Theorem, 2004]\label{thm:GMT}
Every infinite sequence of graphs $G_1, G_2, \ldots$ contains two graphs
$G_i \preceq G_j$ with $i < j$. In particular, any minor-closed family of
graphs has a finite set of forbidden minors.
\end{theorem}

\begin{theorem}[Structure Theorem for $K_k$-minor-free graphs, Robertson--Seymour]\label{thm:struct}
Graphs with no $K_k$ minor have a tree-like structure (bounded treewidth
after surface embeddings and apex vertex removal). More precisely, they can
be built from graphs embedded on surfaces of bounded genus, with a bounded
number of apex vertices and vortices.
\end{theorem}

%% =============================================================================
\section{Verified Cases: \texorpdfstring{$k \leq 4$}{k <= 4}}
%% =============================================================================

\begin{theorem}[Hadwiger's conjecture for $k \leq 3$]\label{thm:k3}
\begin{enumerate}[label=(\roman*)]
  \item $k = 1$: Every nonempty graph has $\chi(G) \geq 1$ and contains $K_1$.
  \item $k = 2$: $\chi(G) \geq 2$ implies $G$ has an edge, hence $K_2 \preceq G$.
  \item $k = 3$: $\chi(G) \geq 3$ implies $G$ contains an odd cycle, hence $K_3 \preceq G$.
\end{enumerate}
\end{theorem}

\begin{proof}
For $k = 3$: Since $\chi(G) \geq 3$, $G$ is not bipartite, so $G$ contains
an odd cycle $C_{2m+1}$ for some $m \geq 1$. Any odd cycle contains $K_3$ as
a minor (contract all but three vertices appropriately), so $K_3 \preceq G$.
\end{proof}

\begin{theorem}[Hadwiger for $k = 4$, Dirac 1952]\label{thm:k4}
Every 4-chromatic graph contains $K_4$ as a minor.
\end{theorem}

\begin{proof}[Proof sketch]
Dirac~\cite{Dirac52} proved that every $4$-chromatic graph contains $K_4$
as a topological minor (i.e., a subdivision of $K_4$). Since a topological
minor is also a minor, the result follows. The proof uses the fact that
$4$-chromatic graphs cannot be planar (by the planar $5$-color theorem) and
a structural analysis via Kempe chains.
\end{proof}

%% =============================================================================
\section{Wagner's Theorem and the Case \texorpdfstring{$k = 5$}{k=5}}
%% =============================================================================

The case $k = 5$ is particularly significant because of its deep connection to
planarity and the Four-Color Theorem.

\begin{theorem}[Wagner, 1937]\label{thm:wagner}
A graph $G$ is planar if and only if it contains neither $K_5$ nor $K_{3,3}$
as a minor.
\end{theorem}

\begin{remark}
This is a minor-version of Kuratowski's theorem (1930), which uses topological
minors (subdivisions) instead of minors. Wagner's theorem is slightly stronger
since minor is a weaker condition than topological minor.
\end{remark}

\begin{theorem}[Four-Color Theorem, Appel--Haken 1977; Robertson et al.\ 1997]\label{thm:4CT}
Every planar graph is $4$-colorable: $\chi(G) \leq 4$ for all planar $G$.
\end{theorem}

\begin{theorem}[Hadwiger for $k = 5$ $\Leftrightarrow$ 4CT, Wagner 1937]\label{thm:k5}
Hadwiger's conjecture for $k = 5$ is equivalent to the Four-Color Theorem.
\end{theorem}

\begin{proof}
($\Leftarrow$) Assume the Four-Color Theorem. Let $G$ be a graph with $\chi(G) \geq 5$.
We must show $G \succeq K_5$.

Since $\chi(G) \geq 5$, by the 4CT $G$ is not planar. By Wagner's theorem
(Theorem~\ref{thm:wagner}), $G$ contains $K_5$ or $K_{3,3}$ as a minor.
If $G \succeq K_5$, we are done. If $G \succeq K_{3,3}$, then since
$\chi(K_{3,3}) = 2$... but wait, $G$ is $(5)$-chromatic. We need a more
careful argument.

The full proof uses Wagner's structure theorem: a graph with no $K_5$ minor
is built from planar graphs and a specific graph $V_8$ (the Wagner graph) via
$3$-sum operations. Such graphs are $4$-colorable by the Four-Color Theorem
applied piecewise. Hence if $\chi(G) \geq 5$, then $G$ must have a $K_5$ minor.

($\Rightarrow$) If Hadwiger holds for $k=5$, then every planar graph $G$
(which has no $K_5$ minor by Wagner's theorem) satisfies $\chi(G) \leq 4$.
\end{proof}

%% =============================================================================
\section{The Case \texorpdfstring{$k = 6$}{k=6}: Robertson--Seymour--Thomas}
%% =============================================================================

\begin{theorem}[Robertson--Seymour--Thomas, 1993]\label{thm:k6}
Every 6-chromatic graph contains $K_6$ as a minor. Equivalently, every graph
with no $K_6$ minor is $5$-colorable.
\end{theorem}

\begin{proof}[Proof outline]
Robertson, Seymour, and Thomas~\cite{RST93} proved that every minimal
counterexample to Hadwiger's conjecture for $k = 6$ would be a graph $G$ that:
\begin{enumerate}[label=(\roman*)]
  \item is $6$-chromatic (hence $\chi(G) = 6$),
  \item contains no $K_6$ minor,
  \item is vertex-transitive (by minimality arguments).
\end{enumerate}

They then showed that any such graph must have an \emph{apex vertex} $v$ such
that $G - v$ is planar. Since $G - v$ is planar, by the Four-Color Theorem
$\chi(G - v) \leq 4$. Coloring $v$ with a fifth color gives $\chi(G) \leq 5$,
a contradiction with $\chi(G) = 6$.

More precisely: if $G$ has no $K_6$ minor, they prove $G$ can be built from
graphs that are either planar or obtained from a planar graph by adding an apex
vertex. This structural decomposition, combined with the Four-Color Theorem
for the planar parts, yields $\chi(G) \leq 5$.
\end{proof}

\begin{corollary}
The cases $k \leq 6$ of Hadwiger's conjecture are theorems. Any proof of
the cases $k \leq 6$ uses the Four-Color Theorem for $k = 5, 6$.
\end{corollary}

%% =============================================================================
\section{Minor-Closed Graph Families and Wagner's Formula}
%% =============================================================================

\begin{theorem}[Hadwiger number and average degree]\label{thm:avgdeg}
If $G$ has average degree $\bar{d}(G) = 2|E|/|V|$, then
\[
  \eta(G) \geq \frac{\bar{d}(G)}{2\sqrt{\log_2 \bar{d}(G)}}.
\]
\end{theorem}

\begin{theorem}[Kostochka, 1984; Thomason, 1984]\label{thm:kostochka}
There exists a constant $c > 0$ such that every graph $G$ with average degree
at least $c \cdot k \sqrt{\log k}$ contains $K_k$ as a minor. The optimal
constant satisfies $c = (\alpha + o(1)) / \sqrt{\ln k}$ where $\alpha = 0.319\ldots$.
\end{theorem}

\begin{corollary}
If Hadwiger's conjecture holds, then any $k$-chromatic graph has average degree
$\Omega(k \sqrt{\log k})$.
\end{corollary}

\begin{remark}
This bound is nearly tight: the complete graph $K_k$ has average degree $k-1$
and chromatic number $k$, so the gap between $k$ and $k \sqrt{\log k}$ shows
there is some slack. However, no example of a $k$-chromatic graph without
$K_k$ minor is known.
\end{remark}

\begin{theorem}[Reed--Seymour, 1998]\label{thm:reedsey}
Every $K_k$-minor-free graph has a proper $(k-1)$-coloring that can be found
in polynomial time, assuming $k \leq 6$ (proven) or Hadwiger's conjecture (open).
\end{theorem}

%% =============================================================================
\section{Kühn--Osthus and Large Chromatic Number}
%% =============================================================================

\begin{theorem}[Kühn--Osthus, 2005]\label{thm:KO}
For every $\varepsilon > 0$, there exists $k_0$ such that for all $k \geq k_0$:
if $G$ has chromatic number $\chi(G) \geq k$, then $G$ contains $K_{\lceil k/(1+\varepsilon)\rceil}$
as a minor. That is:
\[
  \eta(G) \geq \frac{\chi(G)}{1 + \varepsilon} \quad \text{for all sufficiently large } \chi(G).
\]
\end{theorem}

\begin{proof}[Proof idea]
Kühn and Osthus~\cite{KO05} use the following approach:
\begin{enumerate}[label=(\roman*)]
  \item Every $k$-chromatic graph contains a $k$-critical subgraph $G'$
        (removing any vertex decreases the chromatic number).
  \item In a $k$-critical graph, $\delta(G') \geq k-1$ (minimum degree).
  \item Using the Thomason bound (Theorem~\ref{thm:kostochka}), high minimum
        degree implies large Hadwiger number.
  \item The $(1+\varepsilon)$ factor accounts for the gap between $\delta$ and
        $\bar{d}$ in the Thomason bound.
\end{enumerate}
\end{proof}

\begin{theorem}[Norin--Song, 2023]\label{thm:NS23}
For graphs $G$ with $\chi(G) \geq k$,
\[
  \eta(G) \geq k^{1 - o(1)}.
\]
Specifically, $\eta(G) \geq k / (\log k)^{0.999}$ for all sufficiently large $k$.
\end{theorem}

This is the current best asymptotic result, but still falls short of the linear
bound $\eta(G) \geq \chi(G) = k$ required by Hadwiger's conjecture.

%% =============================================================================
\section{Fractional and Linear Hadwiger}
%% =============================================================================

\begin{definition}
The \emph{fractional chromatic number} $\chi_f(G)$ is the infimum of $p/q$
over all proper colorings of $G^q$ (the $q$-fold lexicographic power of $G$)
that use $p$ colors. It satisfies $\chi_f(G) \leq \chi(G)$.
\end{definition}

\begin{theorem}[Fractional Hadwiger, Reed--Seymour 1998]\label{thm:fracHad}
$\chi_f(G) \leq 2\eta(G)$ for all graphs $G$.
\end{theorem}

\begin{conjecture}[Linear Hadwiger]\label{conj:linHad}
There exists an absolute constant $c$ such that $\chi(G) \leq c \cdot \eta(G)$
for all graphs $G$.
\end{conjecture}

\begin{remark}
Hadwiger's conjecture (Conjecture~\ref{conj:hadwiger}) is $\chi(G) \leq \eta(G)$
(constant $c = 1$). The Linear Hadwiger conjecture asks for any fixed $c$,
and would already be a major step. The current best bound from Theorem~\ref{thm:NS23}
gives $\chi(G) \leq (\log k)^{0.999} \cdot \eta(G)$ for large $k$, still logarithmic.
\end{remark}

%% =============================================================================
\section{The Odd Hadwiger Conjecture}
%% =============================================================================

\begin{conjecture}[Odd Hadwiger, Gerards--Seymour]\label{conj:odd}
If $G$ has no $K_k$ as an \emph{odd minor} (i.e., a minor with all branch sets
of odd size or some specific parity condition), then $\chi(G) \leq 2(k-1)$.
\end{conjecture}

\begin{theorem}[Geelen, Gerards, Reed, Seymour, Vetta, 2009]\label{thm:oddHad}
Odd Hadwiger's conjecture holds for $k \leq 4$.
\end{theorem}

%% =============================================================================
\section{Degeneracy and Coloring in Minor-Free Graphs}
%% =============================================================================

\begin{definition}
A graph $G$ is \emph{$d$-degenerate} if every subgraph has a vertex of degree
$\leq d$. Every $d$-degenerate graph is $(d+1)$-colorable.
\end{definition}

\begin{theorem}[Mader, 1967]\label{thm:mader}
Every graph with no $K_k$ minor has average degree less than $2^{k-2}$. In
particular, every $K_k$-minor-free graph is $(2^{k-2})$-degenerate, hence
$(2^{k-2}+1)$-colorable.
\end{theorem}

\begin{remark}
For $k = 5$: Mader's bound gives $5$-colorability (since $2^{4-2}+1 = 5$),
without the Four-Color Theorem. For $k = 6$: $9$-colorability without further
work, but the RST theorem improves this to $5$. The exponential gap between
Mader's bound $2^{k-2}$ and Hadwiger's prediction $k-1$ is enormous for large $k$.
\end{remark}

%% =============================================================================
\section{Current Status and Open Problems}
%% =============================================================================

\begin{conjecture}[Hadwiger for $k \geq 7$ (Open)]\label{conj:open7}
For $k \geq 7$: every graph with $\chi(G) \geq k$ contains $K_k$ as a minor.
\end{conjecture}

The current state of knowledge:
\begin{center}
\begin{tabular}{cll}
\toprule
$k$ & Status & Reference \\
\midrule
$1, 2, 3$ & Trivial & Classical \\
$4$ & Proved & Dirac (1952) \\
$5$ & Proved $\Leftrightarrow$ 4CT & Wagner (1937) + Appel-Haken (1977) \\
$6$ & Proved (uses 4CT) & Robertson-Seymour-Thomas (1993) \\
$7$ & \textbf{Open} & — \\
$\geq 7$ & \textbf{Open} & — \\
$\geq k_0$ (large) & $\eta(G) \geq k/(\log k)^{0.999}$ & Norin-Song (2023) \\
\bottomrule
\end{tabular}
\end{center}

\medskip

The primary open problems are:
\begin{enumerate}[label=(\arabic*)]
  \item Prove or disprove Hadwiger's conjecture for $k = 7$.
  \item Is there an proof for all $k$ that does not use the Four-Color Theorem?
  \item Find an explicit graph with $\chi(G) = k$ but $\eta(G) < k$ (counterexample)
        or prove no such graph exists.
  \item Improve the asymptotic bound: show $\eta(G) \geq c \cdot \chi(G)$ for
        some $c > 1/(\log \chi(G))$.
  \item Prove the Linear Hadwiger conjecture.
\end{enumerate}

%% =============================================================================
\section{Conclusion}
%% =============================================================================

Hadwiger's conjecture unifies two central themes of graph theory: chromatic numbers
and graph minors. Its proof for $k \leq 6$ already required one of the deepest
theorems in mathematics (the Four-Color Theorem), and the cases $k \geq 7$ are
expected to be dramatically harder.

The Robertson--Seymour graph minor theory provides powerful structural tools,
but even these do not suffice to resolve the conjecture for $k \geq 7$. The
problem sits at the intersection of extremal graph theory, topological graph
theory, and computational complexity, and its resolution would likely require
fundamentally new ideas.

\begin{thebibliography}{99}

\bibitem{Hadwiger1943}
H.~Hadwiger,
\textit{Über eine Klassifikation der Streckenkomplexe},
Vierteljschr.\ Naturforsch.\ Ges.\ Zürich \textbf{88} (1943), 133--142.

\bibitem{Wagner37}
K.~Wagner,
\textit{Über eine Eigenschaft der ebenen Komplexe},
Math.\ Ann.\ \textbf{114} (1937), 570--590.

\bibitem{Dirac52}
G.~A.\ Dirac,
\textit{A property of $4$-chromatic graphs and some remarks on critical graphs},
J.\ London Math.\ Soc.\ \textbf{27} (1952), 85--92.

\bibitem{RST93}
N.~Robertson, P.~D.\ Seymour, R.~Thomas,
\textit{Hadwiger's conjecture for $K_6$-free graphs},
Combinatorica \textbf{13} (1993), no.\ 3, 279--361.

\bibitem{AppHak77}
K.~Appel, W.~Haken,
\textit{Every planar map is four colorable},
Bull.\ Amer.\ Math.\ Soc.\ \textbf{82} (1976), 711--712.

\bibitem{RSST97}
N.~Robertson, D.~Sanders, P.~D.\ Seymour, R.~Thomas,
\textit{The four-colour theorem},
J.\ Combin.\ Theory Ser.\ B \textbf{70} (1997), no.\ 1, 2--44.

\bibitem{Kostochka84}
A.~V.\ Kostochka,
\textit{Lower bound of the Hadwiger number of graphs by their average degree},
Combinatorica \textbf{4} (1984), no.\ 4, 307--316.

\bibitem{Thomason84}
A.~Thomason,
\textit{An extremal function for contractions of graphs},
Math.\ Proc.\ Cambridge Philos.\ Soc.\ \textbf{95} (1984), 261--265.

\bibitem{KO05}
D.~Kühn, D.~Osthus,
\textit{Forcing unbalanced complete bipartite minors},
European J.\ Combin.\ \textbf{26} (2005), 75--81.

\bibitem{ReedSey98}
B.~Reed, P.~D.\ Seymour,
\textit{Fractional colouring and Hadwiger's conjecture},
J.\ Combin.\ Theory Ser.\ B \textbf{74} (1998), 147--152.

\bibitem{NorinSong23}
S.~Norin, L.~Song,
\textit{Breaking the degeneracy barrier for coloring graphs with no $K_t$ minor},
Adv.\ Math.\ \textbf{422} (2023), Paper no.\ 109020.

\bibitem{Mader67}
W.~Mader,
\textit{Homomorphieeigenschaften und mittlere Kantendichte von Graphen},
Math.\ Ann.\ \textbf{174} (1967), 265--268.

\end{thebibliography}

\end{document}
